\documentclass[11pt,a4paper]{article}
\usepackage[textwidth=17cm, top=2.8cm, bottom=2.8cm, hmarginratio=1:1, headheight=13.6pt]{geometry}
\usepackage[T1]{fontenc}
\IfFileExists{atkinson.sty}{\usepackage[sfdefault]{atkinson}}{}  % Atkinson Hyperlegible, if available
\usepackage{microtype}            % also provides \textls (letter spacing)
\usepackage{amsmath}
\usepackage{amssymb}
\usepackage{graphicx}
\usepackage{xcolor}
\usepackage{booktabs}
\usepackage{array}
\usepackage{upquote}              % straight quotes in \verb and in listings
\usepackage{float}                % figures with the placement H stay where they are
\usepackage{minted}
\usepackage{fancyhdr}
\usepackage{lastpage}
\usepackage{hyperref}
\hypersetup{colorlinks = true, linkcolor = blue, urlcolor = blue}

\setlength{\emergencystretch}{2em}
\setlength{\parindent}{0pt}
\setlength{\parskip}{0.5\baselineskip}

\DeclareUnicodeCharacter{21D2}{\ensuremath{\Rightarrow}}  % used in the output files of Larkita

\definecolor{sepia}{rgb}{0.95, 0.95, 0.9}
\newcommand{\blue}[1]{{\color{blue}#1}}

\pagestyle{fancy}
\fancyhf{}
\fancyhead[L]{\rightmark}
\fancyhead[R]{A Tutorial Introduction to \Larkita}
\fancyfoot[C]{--- \thepage/\pageref{LastPage}\ ---}
\renewcommand{\sectionmark}[1]{\markright{\thesection. \ #1}}
\fancypagestyle{plain}{\fancyhf{}\fancyfoot[C]{--- \thepage/\pageref{LastPage}\ ---}%
                       \renewcommand{\headrulewidth}{0pt}}

% \simtextsc{Larkita}: the first letter as a capital, the rest as shrunken capitals.
% The font has no small capitals, hence they are faked.
\newcommand{\fakesc}[1]{\resizebox{!}{0.99ex}{\textls[90]{\MakeUppercase{#1}}}}
\makeatletter
\def\@simtextsc#1#2\@nil{\MakeUppercase{#1}\fakesc{#2}}
\newcommand{\simtextsc}[1]{%
  \ifmmode\text{\@simtextsc#1\@nil}\else\@simtextsc#1\@nil\fi}
\makeatother
\newcommand{\Larkita}{\simtextsc{Larkita}}
\newcommand{\Lark}{\simtextsc{Lark}}

% Every listing is read from a file in the directory examples, hence the listings show the
% actual code.  The outputs in examples/output are produced by examples/run_examples.py.
% Grammars and programs are set as figures with the label fig:<file name>.  They do not float, so every
% program is followed by its output.  The program while.py is too long for one page, so it is split into
% two figures.  Outputs are put into the environment unbreakable,
% so they are not split across pages.
\setminted{frame = lines, framesep = 0.3cm, bgcolor = sepia, tabsize = 4}
\newenvironment{unbreakable}{\par\noindent\begin{minipage}{\linewidth}}{\end{minipage}\par}
% \codefigure[minted options]{language}{file}{label}{caption}
\newcommand{\codefigure}[5][]{\begin{figure}[H]%
  \inputminted[numbers = left, numbersep = 0.3cm, #1]{#2}{examples/#3}%
  \vspace*{-0.3cm}%
  \caption{#5}\label{#4}\end{figure}}
\newcommand{\grammarfile}[2]{\codefigure{text}{#1}{fig:#1}{#2}}
\newcommand{\pythonfile}[2]{\codefigure{python3}{#1}{fig:#1}{#2}}
\newcommand{\resultfile}[2][\normalsize]{\begin{unbreakable}%
  \inputminted[fontsize = #1, breaklines]{text}{examples/output/#2.txt}\end{unbreakable}}

\title{A Tutorial Introduction to \Larkita}
\author{Karl Stroetmann}
\date{\today}

\begin{document}
\maketitle
\thispagestyle{plain}

\begin{abstract}
  \Larkita\ is a parser generator for \textsl{Python}.  It reads a context-free grammar and produces a
  canonical LR(1) parser that returns the abstract syntax tree of its input as a \blue{nested tuple}.  The
  shape of these tuples follows from the grammar, so there is no tree-building code to write.  Furthermore,
  \Larkita\ describes the LR(1) automaton it has generated in an \blue{output file}, which makes conflicts easy
  to understand.  If the grammar is LALR(1), \Larkita\ can also produce the much smaller LALR(1) parser: On
  request, it merges states that only differ in their look-ahead tokens, as long as this is safe.  This
  tutorial explains how to write grammars for \Larkita, how the abstract syntax trees are
  built, how to process them with structural pattern matching, and how to deal with syntax errors and
  conflicts.
\end{abstract}

\tableofcontents

\section{Introduction}
\Larkita\ is a light-weight version of the parser generator
\href{https://github.com/lark-parser/lark}{\Lark}.  It differs from \Lark\ in three ways:
\begin{enumerate}
\item By default, it generates a \blue{canonical LR(1) parser}.  These parsers can handle more grammars than
      LALR(1) parsers, and when a grammar has a conflict, the LR(1) states show directly where the conflict
      comes from.  With the option \verb|merge_states=True|, \Larkita\ merges states that only differ in
      their look-ahead tokens, as long as this causes no new conflicts and does not change how the input is
      split into tokens.  If the grammar is LALR(1) and the basic scanner is used, the result is the LALR(1)
      parser, which has far fewer states than an LR(1) parser (see Section \ref{sec:merging}).
\item The parser returns the abstract syntax tree as a \blue{nested tuple}.  For example, with the grammar
      shown in Figure \ref{fig:expr.g}, the string \verb|"x := a + 3;"| is turned into the tuple
      \\[0.2cm]
      \hspace*{0.3cm}
      \verb|(':=', ('IDENTIFIER', 'x'),|
      \\
      \hspace*{2.0cm}
      \verb|('+', ('IDENTIFIER', 'a'), ('NUMBER', '3')))|.
      \\[0.2cm]
      Tuples like these are easy to print and easy to process with the \verb|match| statement of
      \textsl{Python}.
\item Every time a parser is generated, \Larkita\ writes an \blue{output file} that lists all LR(1) states.
      Every item of a state is annotated with its look-ahead tokens.
\end{enumerate}
Grammars are written in the EBNF notation of \Lark, with a few additions that control the shape of the
abstract syntax trees.

\subsection{Installation}
\Larkita\ is written in pure \textsl{Python} and needs no other packages.  It can be installed directly from
\href{https://github.com/karlstroetmann/Larkita}{GitHub}:
\begin{unbreakable}
\begin{minted}{text}
pip install git+https://github.com/karlstroetmann/Larkita.git
\end{minted}
\end{unbreakable}
\Larkita\ runs with \textsl{Python} 3.8 or newer.  The examples in this tutorial use the \verb|match|
statement and therefore need \textsl{Python} 3.10 or newer.

\subsection{The examples}
All grammars and programs shown in this tutorial are contained in the directory \texttt{tutorial/examples}
of the \Larkita\ repository.  The listings in this document are read from these files and the outputs shown
are the outputs that the programs actually produce: The script \texttt{run\_examples.py} runs all examples
and stores their outputs in the directory \texttt{examples/output}, from where they are included into this
document.

\section{A First Parser}
The file \texttt{expr.g} shown in Figure \ref{fig:expr.g} contains a grammar for assignments and arithmetic
expressions.
\grammarfile{expr.g}{The grammar \texttt{expr.g} for assignments and arithmetic expressions.}
\begin{itemize}
\item Line 1 is a comment.  Comments start with \verb|//| and extend to the end of the line.
\item A \blue{rule} consists of the name of a syntactical variable, a colon, and a list of
      \blue{alternatives} that are separated by the character ``\verb+|+''.  The names of syntactical variables
      are written in lower case.  The first rule defines the \blue{start symbol} of the grammar, which is
      \texttt{stmnt} in this case.
\item Strings enclosed in double quotes like \verb|":="| or \verb|"+"| are \blue{literals}: They stand for
      exactly the characters written between the quotes.
\item Names written in upper case like \texttt{IDENTIFIER} and \texttt{NUMBER} denote \blue{terminals}.
      These terminals are defined in line 18 -- 20:  The command
      \\[0.2cm]
      \hspace*{1.3cm}
      \verb|%import common.CNAME -> IDENTIFIER|
      \\[0.2cm]
      imports the terminal \texttt{CNAME} from the library \texttt{common.lark}, which is part of \Larkita,
      and renames it into \texttt{IDENTIFIER}.  \texttt{CNAME} matches names as they are used in the
      programming language \texttt{C}.  Similarly, \texttt{NUMBER} matches integers and floating point
      numbers, and \texttt{WS} matches white space.
\item Line 21 tells the scanner to skip white space.
\item The annotations \verb|%ignore ";"| in line 3 and \verb|%ignore "(", ")"| in line 14 are specific to
      \Larkita.  They tell \Larkita\ that the semicolon of an assignment and the parentheses should not
      appear in the abstract syntax tree.  The semicolon in line 4 is kept:  It is the only literal of this
      alternative, and without it the abstract syntax tree of the statement would just be the tree of the
      expression.  We will discuss \verb|%ignore| in Section \ref{sec:ast}.
\end{itemize}
The program \texttt{first.py} in Figure \ref{fig:first.py} generates a parser from this grammar and parses a
string.
\pythonfile{first.py}{The program \texttt{first.py} parses a string with the grammar \texttt{expr.g}.}
\begin{itemize}
\item The class method \verb|Larkita.open| reads the grammar from a file and returns a parser.  If the grammar
      is given as a string, we use the constructor instead:  \verb|Larkita(grammar)|.
\item The method \texttt{parse} takes the string to parse and returns its abstract syntax tree.  We print this
      tree with the function \texttt{pprint}, which inserts line breaks and indentation:
\end{itemize}
\resultfile{first}
Every node of this tree is a tuple.  The first component of a tuple is called its \blue{head}: It contains the
literals of the rule that has been used, except for the ignored ones.  The remaining components are the \blue{children}.  For example, the
tuple \verb|('+', ('IDENTIFIER', 'a'), ('NUMBER', '3'))| represents the expression \verb|a + 3|.  Its head is
the operator \verb|'+'| and its children represent the operands.  A token like \texttt{a} is represented as a
pair consisting of the name of the terminal and the text that has been read.  Section \ref{sec:ast} explains
the construction of the abstract syntax tree in detail.

When \Larkita\ generates a parser, it also writes an output file that describes the parser.  As the grammar
has been read from the file \texttt{expr.g}, this file is called \texttt{expr.out} and is written next to the
grammar.  If the grammar is given as a string, the output file is called \texttt{grammar.out} and is written
into the current directory.  The parameter \texttt{out\_file} sets another path.  Section \ref{sec:output}
explains the output file.

\section{Writing Grammars}
\Larkita\ reads grammars in the EBNF notation of \Lark.  This section summarizes this notation.

\subsection{Rules}
The right hand side of a rule can use the following operators:
\begin{center}
\begin{tabular}{ll}
\toprule
\textbf{Notation}        & \textbf{Meaning} \\
\midrule
\verb+a | b+             & either \texttt{a} or \texttt{b} \\
\verb|(a b)|             & grouping \\
\verb|[a b]|             & optional: \texttt{a b} or nothing \\
\verb|a?|                & optional: \texttt{a} or nothing \\
\verb|a*|                & zero or more repetitions of \texttt{a} \\
\verb|a+|                & one or more repetitions of \texttt{a} \\
\verb|a ~ 3|             & exactly three repetitions of \texttt{a} \\
\verb|a ~ 2..4|          & between two and four repetitions of \texttt{a} \\
\bottomrule
\end{tabular}
\end{center}
\Larkita\ translates these operators into ordinary grammar rules.  The output file lists the resulting rules.

Alternatives are usually written on lines of their own.  An alternative may be empty, so the rule
\begin{unbreakable}
\begin{minted}{text}
args: expr ("," expr)*
    |
\end{minted}
\end{unbreakable}
describes a possibly empty list of expressions.  A rule can be given a \blue{priority} by appending a dot and
a number to its name, as in \verb|expr.2: ...|.  Priorities are used to resolve reduce/reduce conflicts, see
Section \ref{sec:priorities}.

\subsection{Terminals}
There are three ways to specify the tokens of a language:
\begin{enumerate}
\item A \blue{literal} is a string enclosed in double quotes, e.g.~\verb|"while"|.  If the letter \texttt{i}
      is appended, as in \verb|"while"i|, upper and lower case are not distinguished.
\item A \blue{regular expression} is enclosed in slashes, e.g.~\verb|/[0-9]+/|.  It can be followed by flags,
      e.g.~\verb|/begin/i|.  The syntax is the syntax of the \textsl{Python} module \texttt{re}.
\item A \blue{terminal} is defined by a rule whose name is written in upper case.  Its right hand side may
      use literals, regular expressions, other terminals, and the operators shown above.  For example:
\begin{unbreakable}
\begin{minted}{text}
DIGIT:  "0".."9"
INT:    DIGIT+
FLOAT:  INT "." INT
STRING: /"[^"]*"/
\end{minted}
\end{unbreakable}
      The notation \verb|"0".."9"| denotes a range of characters.  A terminal can be given a priority, as in
      \verb|FLOAT.2: INT "." INT|.  If two terminals match the input, the scanner prefers the terminal with
      the higher priority.
\end{enumerate}
The library \texttt{common.lark} defines the following terminals, among others:
\begin{center}
\begin{tabular}{l>{\raggedright\arraybackslash}p{10.3cm}}
\toprule
\textbf{Terminal}        & \textbf{Matches} \\
\midrule
\texttt{INT}, \texttt{SIGNED\_INT}             & natural numbers and integers \\
\texttt{FLOAT}, \texttt{NUMBER}                & floating point numbers; integers or floating point numbers \\
\texttt{CNAME}                                 & names, i.e.~a letter or an underscore followed by letters,
                                                 digits, and underscores \\
\texttt{ESCAPED\_STRING}                       & strings in double quotes \\
\texttt{WS}, \texttt{WS\_INLINE}, \texttt{NEWLINE}
                                               & white space; white space without line breaks; line breaks \\
\texttt{CPP\_COMMENT}, \texttt{C\_COMMENT}     & comments of the form \verb|// ...| and \verb|/* ... */| \\
\texttt{SH\_COMMENT}                           & comments of the form \verb|# ...| \\
\bottomrule
\end{tabular}
\end{center}
They are imported with \verb|%import common.NAME|; with \verb|%import common.NAME -> NEW_NAME| the terminal
is renamed.

\subsection{The scanner}
\label{sec:scanner}
The scanner splits the input into tokens.  The command \verb|%ignore NAME| tells the scanner to skip the
tokens of the terminal \texttt{NAME}, e.g.~white space or comments.  The method \texttt{lex} runs only the
scanner, which is useful when a grammar does not work as expected.  Figure \ref{fig:lex.py} shows an example.
\pythonfile{lex.py}{The program \texttt{lex.py} splits a string into tokens.}
The method \texttt{lex} returns an iterator over the tokens.  Every token has a type, a value, a line, and a
column.  The program prints them in the form \texttt{line:column}:
\resultfile{lex}
\Larkita\ invents names for the terminals that are given as literals, e.g.~\texttt{STAR} for \verb|"*"|.  If a
literal has no obvious name, it is called \verb|__ANON_0|, \verb|__ANON_1|, etc.  The output file and the
error messages show literals in quotes instead.

Keywords are handled automatically:  If a literal like \verb|"while"| can also be matched by the terminal of
the identifiers, the string \texttt{while} is recognized as the keyword, but \texttt{whiles} or
\texttt{while1} are still recognized as identifiers.

By default, the scanner is \blue{contextual}:  In every state of the parser, it only tries to match those
terminals that the parser can accept in this state.  This avoids many collisions between terminals.  With the
option \verb|lexer='basic'|, the scanner always tries all terminals.

\section{How the Abstract Syntax Tree is Built}
\label{sec:ast}
This section explains how \Larkita\ builds the abstract syntax tree.  Every example shows a small grammar,
followed by some strings and their abstract syntax trees.  The examples are produced by the program
\texttt{shapes.py}, which adds the following terminal definitions to every grammar:
\begin{unbreakable}
\begin{minted}{text}
%import common.CNAME -> NAME
%import common.INT -> NUM
%import common.WS
%ignore WS
\end{minted}
\end{unbreakable}

\subsection{The head}
When the parser reduces with an alternative, the string literals of this alternative form the
\blue{head} of a new tuple.  All other symbols of the alternative provide the \blue{children} of the
tuple:  A syntactical variable contributes its abstract syntax tree, while a named terminal contributes a pair
of the form \verb|(NAME, text)|, where \texttt{text} is the text that has been read.  This text is never
converted, so the number 1 is represented as \verb|('NUM', '1')|.  If an alternative contains exactly one
literal, the head is this literal:
\resultfile{shapes-head}
If an alternative contains several literals, the head is the tuple of these literals in the order in which
they occur in the alternative:
\resultfile{shapes-heads}

\subsection{Pass-through}
If an alternative has no head and consists of a single symbol, then no new tuple is built.  Instead, the
abstract syntax tree of this symbol is returned.  This rule is called \blue{pass-through}.  It guarantees that
alternatives like \verb|sum: atom| leave no trace in the abstract syntax tree:
\resultfile{shapes-pass-through}
The string \texttt{42} is parsed with the alternatives \verb|sum: atom| and \verb|atom: NUM|, but the result is
just the pair \verb|('NUM', '42')|.

\subsection{The rule name as head}
If an alternative has no head and is not a pass-through, the name of the rule is used as head:
\resultfile{shapes-rule-name}
Pass-through is decided by the grammar, not by the input:  The alternative \verb|program: stmnt*|
always builds a tuple, even if the input contains only one statement.  For the empty string, it builds the
tuple \verb|('program',)|, which has no children.

\subsection{Repetitions and groups}
The children produced by repetitions and groups are inserted into the tuple of the alternative in which they
occur.  Literals that occur inside repetitions, optional parts, and groups belong to the head:
\resultfile{shapes-flatten}
Note that the head is determined by the grammar, not by the input:  The head of the abstract syntax tree of
\texttt{[1]} contains a comma, although no comma has been read.  Therefore, all lists produced by this
alternative have the same head.

\subsection{Optional parts}
If an optional part is missing, it is simply left out:
\resultfile{shapes-optional}
There is no placeholder like \texttt{None}, so the number of children tells whether the optional part has
been present.

\subsection{Ignoring literals}
Some literals are needed for parsing, but they carry no information once the abstract syntax tree has been
built.  For example, parentheses determine the structure of an expression, but this structure is already
represented by the tuples.  Writing \verb|%ignore| followed by a list of literals at the end of an
alternative removes these literals from the head of this alternative.  In the following example, the
alternative \verb|atom: "(" sum ")"| therefore has no head and is a pass-through:
\resultfile{shapes-ignore}
If some literals should be removed from all alternatives of the grammar, the command \verb|%ignore:| (with a
colon) is used.  In the following example, the alternative has no head left.  As it is not a pass-through, the
name of the rule becomes the head:
\resultfile{shapes-ignore-global}
Note that there are three different forms of \verb|%ignore|:
\begin{center}
\begin{tabular}{l>{\raggedright\arraybackslash}p{10.3cm}}
\toprule
\textbf{Form}                    & \textbf{Effect} \\
\midrule
\verb|%ignore WS|                & the scanner skips the tokens of the terminal \texttt{WS} \\
\verb|a: ... %ignore "(", ")"|   & the literals are removed from the head of this alternative \\
\verb|%ignore: "(", ")"|         & the literals are removed from the heads of all alternatives \\
\bottomrule
\end{tabular}
\end{center}

\subsection{Aliases}
An \blue{alias} replaces the head of an alternative.  It is written at the end of the alternative, either as
a name or as a string:  \verb|-> call| and \verb|-> "call"| are equivalent.
\resultfile{shapes-alias}
Aliases are useful when the same literals are used for different purposes.  Here, the unary minus and the
binary minus both use the literal \verb|"-"|.  Without the alias, both would produce the head \verb|'-'|.
Although the number of children would still tell them apart, the alias \texttt{neg} makes the abstract syntax
tree easier to read.

\subsection{Summary}
The following list summarizes the rules.  Every alternative is followed by the abstract syntax tree that it
produces, where \texttt{ast(x)} denotes the abstract syntax tree of \texttt{x}:
\begin{unbreakable}
\begin{minted}{text}
expr: expr "+" product
    ⇒ ('+', ast(expr), ast(product))
stmnt: IDENTIFIER ":=" expr ";"
    ⇒ ((':=', ';'), ('IDENTIFIER', name), ast(expr))
atom: IDENTIFIER
    ⇒ ('IDENTIFIER', name)
expr: product
    ⇒ ast(product)
factor: "(" expr ")"    %ignore "(", ")"
    ⇒ ast(expr)
atom: NAME "(" args ")"    -> "call"
    ⇒ ('call', ('NAME', name), ast(args))
args: expr ("," expr)*
    ⇒ (',', ast(expr_1), ..., ast(expr_n))
program: stmnt*
    ⇒ ('program', ast(stmnt_1), ..., ast(stmnt_n))
ret: "return" [expr] ";"
    ⇒ (('return', ';'),)  or  (('return', ';'), ast(expr))
\end{minted}
\end{unbreakable}

\section{Processing the Abstract Syntax Tree}
Abstract syntax trees that are represented as nested tuples are best processed with the \verb|match|
statement of \textsl{Python}.  As an example, we implement a simple calculator.  Its grammar
\texttt{calc.g}, which is shown in Figure \ref{fig:calc.g}, extends the grammar \texttt{expr.g}:  A program
is a sequence of statements, the statement \verb|print expr;| prints the value of an expression, and there is
a unary minus.
\grammarfile{calc.g}{The grammar \texttt{calc.g} of a simple calculator.}
The program \texttt{calc.py} in Figure \ref{fig:calc.py} parses a small program, prints its abstract syntax
tree, and executes it.
\pythonfile{calc.py}{The program \texttt{calc.py} implements the calculator.}
\begin{itemize}
\item The function \texttt{evaluate} computes the value of an expression.  Every \verb|case| corresponds to an
      alternative of the grammar:  The pattern \verb|('+', lhs, rhs)| matches the tuples produced by the
      alternative \verb|expr: expr "+" product| and binds the variables \texttt{lhs} and \texttt{rhs} to the
      abstract syntax trees of the operands.  The alias \texttt{neg} distinguishes the unary minus from the
      binary minus.
\item The function \texttt{execute} executes a program.  As the head of the abstract syntax tree of a program
      is the rule name \texttt{program}, the statements are the components \verb|program[1:]|.  As the
      semicolons are ignored, the head of an assignment is \verb|':='| and the head of a \verb|print|
      statement is \verb|'print'|.
\end{itemize}
The program prints the following:
\resultfile{calc}

\section{A Small Interpreter}
\label{sec:interpreter}
Our next example is an interpreter for a small imperative language with assignments, \verb|print|
statements, \verb|if| statements, \verb|while| loops, and blocks.  Its grammar \texttt{while.g} in Figure
\ref{fig:while.g} shows how \verb|%ignore:| and aliases keep the abstract syntax tree small.
\grammarfile{while.g}{The grammar \texttt{while.g} of a small imperative language.}
Instead of ignoring literals in every alternative separately, line 29 removes parentheses, braces, and
semicolons from all heads.  Hence:
\begin{itemize}
\item An assignment \verb|x := e;| is represented as \verb|(':=', ('IDENTIFIER', 'x'), ast(e))|.
\item A statement \verb|print(e);| is represented as \verb|('print', ast(e))|.
\item The head of an \verb|if| statement would be \verb|('if', 'else')|.  The alias in line 7 replaces it with
      \verb|'if'|, so an \verb|if| statement is represented as \verb|('if', ast(c), ast(s1), ast(s2))|.
\item A loop \verb|while (c) s| is represented as \verb|('while', ast(c), ast(s))|.
\item A block has no head left.  Without the alias in line 9, the rule name \texttt{stmnt} would be used as
      head.  With the alias, a block is represented as \verb|('block', ast(s1), ..., ast(sn))|.
\item The alternative \verb|factor: "(" expr ")"| in line 25 has no head and is therefore a pass-through.
\end{itemize}
Line 33 and 35 tell the scanner to skip comments.  Figure \ref{fig:while.py:1} shows the interpreter and
Figure \ref{fig:while.py:2} shows two programs that are executed by it.
\codefigure[lastline = 40]{python3}{while.py}{fig:while.py:1}{The interpreter \texttt{while.py}, part 1:
  the evaluation of expressions and the execution of statements.}
\codefigure[firstline = 42, firstnumber = 42]{python3}{while.py}{fig:while.py:2}{The interpreter
  \texttt{while.py}, part 2: two programs are parsed and executed.}
\begin{itemize}
\item As all binary operators produce tuples of the form \verb|(op, lhs, rhs)|, the function
      \texttt{evaluate} needs only one case for them.  The dictionary \texttt{OPERATORS} maps every operator to
      the corresponding function of the module \texttt{operator}.
\item In the function \texttt{execute}, the pattern \verb+('program' | 'block', *body)+ matches both programs
      and blocks.  The star collects the statements of the program or the block in the list \texttt{body}.
\end{itemize}
The program prints the abstract syntax tree of the first program, followed by the outputs of both programs:
\resultfile{while}

\section{Syntax Errors}
If the input contains a syntax error, the method \texttt{parse} raises an exception.  All exceptions of
\Larkita\ are subclasses of \texttt{LarkitaError}.  There are two kinds of syntax errors:
\begin{itemize}
\item \texttt{UnexpectedCharacters} is raised by the scanner if none of the terminals matches the input.
\item \texttt{UnexpectedToken} is raised by the parser if it reads a token that it cannot accept in its
      current state.  The attribute \texttt{token} contains this token and the attribute \texttt{expected}
      contains the set of terminals that the parser would have accepted.  If the input ends prematurely, the
      token has the type \verb|$END|.
\end{itemize}
Both exceptions are subclasses of \texttt{UnexpectedInput}.  They provide the attributes \texttt{line} and
\texttt{column}, which give the position of the error, and the method \verb|get_context(text)|, which returns
the line containing the error together with a marker that points to the error.  The program
\texttt{errors.py} in Figure \ref{fig:errors.py} checks three erroneous programs for the grammar
\texttt{while.g}.  It prints the following:
\pythonfile{errors.py}{The program \texttt{errors.py} reports the positions of syntax errors.}
\resultfile{errors}
In the last example, the semicolon after \verb|print(x)| is missing.  As the input ends prematurely, the
error is reported at the position of the last token.  The message of an exception describes the error in
more detail.  The program \texttt{message.py} in Figure \ref{fig:message.py} prints such a message:
\pythonfile{message.py}{The program \texttt{message.py} prints the message of a syntax error.}
\resultfile[\small]{message}

\section{The Output File}
\label{sec:output}
When \Larkita\ generates a parser, it writes a description of this parser into an output file.  For the
grammar \texttt{expr.g}, the output file \texttt{expr.out} begins as follows:
\resultfile[\small]{expr-header}
Next, the file lists the rules of the grammar.  Every rule has a number, and every rule is followed by the
abstract syntax tree that a reduction with this rule produces:
\resultfile[\small]{expr-rules}
Rule (0) is added by \Larkita: Its left hand side is the new start symbol \verb|$root_stmnt|.  If the grammar
uses operators like \verb|*| or \verb|?|, the list contains the auxiliary rules that \Larkita\ has introduced
for them.  Then the file lists the conflicts of the grammar, if there are any.  Finally, it describes all
states of the parser.  For example, state 2 is the state that the parser enters after it has read an
identifier at the beginning of a statement:
\resultfile[\small]{expr-state}
\begin{itemize}
\item The first part of a state lists its \blue{items}.  An item is a rule with a dot \verb|•| marking how
      much of the right hand side has already been read.  It is followed by the set of \blue{look-ahead
      tokens}: These are the tokens that can follow the left hand side of the rule when the parser reduces
      with this rule.  Items that have been added by the closure operation are marked with a \verb|+| at the
      beginning of the line.  State 2 has no such items.
\item The second part lists the \blue{actions} of the state.  For every token that the parser can accept in
      this state, it says whether the parser shifts the token or reduces with a rule.  Furthermore, it lists the
      \blue{gotos}, i.e.~the states the parser enters after it has reduced to a syntactical variable.
\end{itemize}
State 2 shows how the parser uses the look-ahead token: After reading an identifier, the parser does not
know yet whether the statement is an assignment or an expression.  If the next token is \verb|":="|, it
shifts this token.  If the next token is an operator or a semicolon, the identifier is part of an expression
and the parser reduces with rule (11).

The parser itself also provides some information about the grammar:  The attributes \texttt{rules},
\texttt{terminals}, \texttt{conflicts}, and \texttt{parse\_table} contain the rules, the terminals, the
conflicts, and the parse table, while \texttt{out\_file} is the path of the output file.

\section{Conflicts}
If the grammar is not LR(1), there are states in which the parser cannot decide which action to take:
These states contain \blue{conflicts}.  A \blue{shift/reduce conflict} occurs when the parser could either
shift the next token or reduce with a rule.  A \blue{reduce/reduce conflict} occurs when it could reduce with
two different rules.

\subsection{Shift/reduce conflicts}
The grammar \texttt{dangling.g} in Figure \ref{fig:dangling.g} describes nested \verb|if| statements.  It
suffers from the \blue{dangling else problem}:  In a statement like \verb|if (a) if (b) x; else y;| it is not
clear whether the \verb|else| belongs to the first or to the second \verb|if|.
\grammarfile{dangling.g}{The grammar \texttt{dangling.g} has a shift/reduce conflict.}
\Larkita\ resolves shift/reduce conflicts in favour of shifting.  In addition, it issues a warning of the
class \texttt{ConflictWarning}.  The program \texttt{dangling.py} in Figure \ref{fig:dangling.py} catches
this warning and prints it:
\pythonfile{dangling.py}{The program \texttt{dangling.py} catches the warning about the conflict.}
\resultfile[\small]{dangling}
Shifting means that an \verb|else| belongs to the innermost \verb|if| that has no \verb|else| yet, which is
the rule used by most programming languages.  The output file \texttt{dangling.out} lists the conflict:
\resultfile[\small]{dangling-conflicts}
The conflict is also shown in the state where it occurs:
\resultfile[\small]{dangling-state}
After reading \verb|if (a) if (b) x;|, the parser could either reduce the inner \verb|if| statement with rule
(1) or shift the token \verb|"else"|.  The look-ahead token \verb|"else"| is contained in the look-ahead sets of
both items.

The warnings are handled by the module \texttt{warnings} of \textsl{Python}.  For example, the statement
\begin{unbreakable}
\begin{minted}{python3}
warnings.simplefilter('error', ConflictWarning)
\end{minted}
\end{unbreakable}
turns these warnings into exceptions, so that every shift/reduce conflict is reported as an error, while
\verb|warnings.simplefilter('ignore', ConflictWarning)| suppresses them.

\subsection{Reduce/reduce conflicts}
\Larkita\ treats reduce/reduce conflicts as errors: It raises a \texttt{GrammarError}.  The grammar
\texttt{look-ahead.g} in Figure \ref{fig:look-ahead.g} is unambiguous, but it is not LR(1).
\grammarfile{look-ahead.g}{The grammar \texttt{look-ahead.g} is unambiguous, but it is not LR(1).}
After reading \verb|"X"|, the parser has to decide whether \verb|"X"| is a \texttt{b} or a \texttt{c}.  This
depends on whether \verb|"U"| is followed by \verb|"V"| or by \verb|"W"|, but the parser can only look at one
token, namely \verb|"U"|.  The program \texttt{look-ahead.py} in Figure \ref{fig:look-ahead.py} tries to
generate a parser for this grammar and prints the message of the exception.  Afterwards, it generates a
parser for the grammar \texttt{look-ahead-solved.g} discussed below and parses two strings:
\pythonfile{look-ahead.py}{The program \texttt{look-ahead.py} reports the reduce/reduce conflict of
  \texttt{look-ahead.g}.}
\resultfile[\small]{look-ahead}
The output file is written even if the grammar has conflicts, so it can be used to find the cause of the
problem:
\resultfile[\small]{look-ahead-state}
This conflict is solved by rewriting the grammar, so that the decision can be postponed until the token
\verb|"U"| has been read.  The result is the grammar \texttt{look-ahead-solved.g} shown in Figure
\ref{fig:look-ahead-solved.g}.
\grammarfile{look-ahead-solved.g}{The grammar \texttt{look-ahead-solved.g} avoids the conflict of
  \texttt{look-ahead.g}.}

\subsection{Priorities}
\label{sec:priorities}
A reduce/reduce conflict can also be resolved by giving one of the rules a higher priority.  In the grammar
\texttt{priority.g} in Figure \ref{fig:priority.g}, a statement consisting of a single name could be both a
declaration of a type without a variable and an expression.  The priority 2 in line 9 resolves this conflict
in favour of the expression.  The program \texttt{priority.py} in Figure \ref{fig:priority.py} parses a
declaration and an expression:
\grammarfile{priority.g}{The grammar \texttt{priority.g} resolves a reduce/reduce conflict with a priority.}
\pythonfile{priority.py}{The program \texttt{priority.py} parses a declaration and an expression.}
\resultfile{priority}
The output file lists the conflict together with its resolution:
\resultfile[\small]{priority-conflicts}
Priorities should be used with care:  The rule with the lower priority is never used to reduce in the
states where the conflict occurs.  Usually, it is better to rewrite the grammar.

\section{Merging States}
\label{sec:merging}
Canonical LR(1) automata can have many states.  For example, the parser for the grammar \texttt{while.g} of
Section \ref{sec:interpreter} has 142 states.  Many of these states have the same \blue{core}, i.e.~they contain the same items
and only differ in the look-ahead sets of these items.  An LALR(1) parser merges all states with the same core,
which yields far fewer states, but this can cause reduce/reduce conflicts that the canonical LR(1) parser does
not have.  With the option \verb|merge_states=True|, \Larkita\ merges states with the same core, as long as
this is safe:
\begin{enumerate}
\item Merging must not cause a new conflict, neither in the merged states nor in their successors.  As two
      merged states have one successor for every symbol, their successors have to be merged, too.
\item Merging must not change how the input is split into tokens.  As explained in Section \ref{sec:scanner}, the
      contextual scanner only tries the terminals that the parser can accept in its current state.  If the
      look-ahead tokens of one state are added to another state, the scanner might read some text as a
      different token.
\end{enumerate}
The program \texttt{merge.py} in Figure \ref{fig:merge.py} generates three parsers for the grammar
\texttt{while.g}: a canonical parser, a parser with merged states, and a parser with merged states that uses
the basic scanner.  It prints the number of states of every parser and parses the string
\verb|"x := 1; else := 2;"|:
\pythonfile{merge.py}{The program \texttt{merge.py} compares parsers with and without merged states.}
\resultfile{merge}
Merging reduces the number of states from 142 to 81.  The output file \texttt{while-merged.out} describes
the merged states.  For example, state 3 has been obtained by merging the LR(1) states 3 and 16.  Its
look-ahead set is the union of their look-ahead sets:
\resultfile[\small]{while-merged-state}
A new section of the output file explains why some states with the same core have been kept apart.  It begins
as follows:
\resultfile[\small]{while-merged}
After the statement \verb|{}|, state 15 expects the next statement, which may start with an identifier, while
state 65 expects the keyword \verb|else|, because the block is the then-branch of an \verb|if| statement.  In
state 15, the contextual scanner reads the string \texttt{else} as an identifier, because the keyword
\verb|"else"| cannot occur here.  If the two states were merged, the scanner would read \texttt{else} as a
keyword in both contexts.  Hence, these states are kept apart.  The output file lists five more pairs of
this kind, one for every other kind of statement.  For example, the parser reaches state 27 after an
assignment like \verb|x := 1;| that is not followed by \verb|else|.  Therefore, both the canonical parser
and the parser with merged states read the string \texttt{else} in \verb|x := 1; else := 2;| as an
identifier.  All other states that have been kept apart only differ in their successors.

The basic scanner always tries all terminals, so merging cannot change the tokens.  Therefore, with the basic
scanner all states with the same core of the grammar \texttt{while.g} are merged and the result is the
LALR(1) automaton with 56 states.  However, the basic scanner always reads \texttt{else} as a keyword, so this
parser rejects \verb|else := 2;|.  The basic scanner without merging would reject it, too.  The output file
\texttt{while-basic.out} reports:
\resultfile[\small]{while-basic}
The grammar \texttt{not-lalr.g} in Figure \ref{fig:not-lalr.g} is LR(1), but it is not LALR(1).
\grammarfile{not-lalr.g}{The grammar \texttt{not-lalr.g} is LR(1), but it is not LALR(1).}
After reading \verb|"v" "x"| and \verb|"w" "x"|, the parser reaches two states with the same core.  In the first
state, \verb|"x"| has to be reduced to \texttt{a} if the next token is \verb|"y"|, and to \texttt{b} if it is
\verb|"z"|.  In the second state, it is the other way round.  Merging these states would therefore cause two
reduce/reduce conflicts, so they are kept apart.  The program \texttt{not-lalr.py} in Figure
\ref{fig:not-lalr.py} generates a parser with merged states and parses four strings:
\pythonfile{not-lalr.py}{The program \texttt{not-lalr.py} generates a parser for \texttt{not-lalr.g} with
  merged states.}
\resultfile{not-lalr}
\resultfile[\small]{not-lalr-merged}
As these are the only states with the same core, merging does not reduce the number of states for this
grammar.

A parser with merged states accepts the same strings and returns the same abstract syntax trees as the
canonical parser, and it detects syntax errors at the same tokens.  However, a merged state may reduce on a
look-ahead token that is an error in some of the states it has been merged from.  Then the parser performs
some reductions before it detects the error, and the list of expected tokens in the error message may
differ.  As the output file of the canonical parser is easier to understand, states are only merged if this
is requested.

\section{Reference}
\subsection{The class \texttt{Larkita}}
\begin{unbreakable}
\begin{minted}{python3}
Larkita(grammar, *, start=None, lexer='contextual', import_paths=(),
        source_path=None, out_file=None, regex=False, g_regex_flags=0,
        merge_states=False)
\end{minted}
\end{unbreakable}
\begin{itemize}
\item \texttt{grammar}: the grammar as a string or as an open file.
\item \texttt{start}: the start symbol.  By default, the first rule defines the start symbol.
\item \texttt{lexer}: either \verb|'contextual'| or \verb|'basic'|.
\item \texttt{import\_paths}: additional directories that are searched for imported grammars.
\item \texttt{source\_path}: the file name of the grammar.  It is used for relative imports and to name the
      output file.
\item \texttt{out\_file}: the path of the output file.
\item \texttt{regex}: if \texttt{True}, the module \texttt{regex} is used instead of the module \texttt{re}.
\item \verb|g_regex_flags|: flags that are applied to all regular expressions, e.g.~\verb|re.IGNORECASE|.
\item \verb|merge_states|: if \texttt{True}, states with the same core are merged, as long as this is safe
      (see Section \ref{sec:merging}).
\end{itemize}
\verb|Larkita.open(filename, rel_to=None, **options)| reads the grammar from a file.  If \verb|rel_to| is
given, the file name is interpreted relative to the directory of the file \verb|rel_to|, so
\verb|Larkita.open('expr.g', rel_to=__file__)| finds a grammar that is stored next to the program.  The other
options are passed to the constructor.

A parser provides the following methods and attributes:
\begin{itemize}
\item \verb|parse(text)| returns the abstract syntax tree of \texttt{text}.
\item \verb|lex(text)| returns an iterator over the tokens of \texttt{text}.
\item \texttt{rules}, \texttt{terminals}, \texttt{conflicts}, \texttt{parse\_table}, and
      \verb|out_file| describe the generated parser.
\end{itemize}

\subsection{Grammar commands}
\begin{itemize}
\item \verb|%import common.NAME| imports a terminal from the library \texttt{common.lark}.
\item \verb|%import .file.NAME| imports a rule or a terminal from the grammar \texttt{file.lark}, which is
      stored in the directory of the current grammar.
\item \verb|-> NEW_NAME| at the end of an import renames the imported symbol.
\item \verb|%ignore NAME| makes the scanner skip the terminal \texttt{NAME}.
\item \verb|%ignore "a", "b"| at the end of an alternative removes these literals from its head, while
      \verb|%ignore: "a", "b"| removes them from all heads.
\item \verb|%declare NAME| declares a terminal that is not defined in the grammar.
\item \verb|%override| and \verb|%extend| replace or extend a rule or a terminal that has been imported.
\end{itemize}

\subsection{Exceptions and warnings}
\begin{center}
\begin{tabular}{l>{\raggedright\arraybackslash}p{11cm}}
\toprule
\textbf{Class}                  & \textbf{Meaning} \\
\midrule
\texttt{LarkitaError}           & the base class of all exceptions of \Larkita \\
\texttt{GrammarError}           & the grammar is erroneous or has an unresolved reduce/reduce conflict \\
\texttt{UnexpectedInput}        & the base class of \texttt{UnexpectedCharacters} and \texttt{UnexpectedToken} \\
\texttt{UnexpectedCharacters}   & no terminal matches the input \\
\texttt{UnexpectedToken}        & the parser cannot accept the next token \\
\texttt{ConflictWarning}        & a shift/reduce conflict has been resolved as shift \\
\bottomrule
\end{tabular}
\end{center}

\subsection{Differences from \Lark}
Readers who know \Lark\ should note the following differences:
\begin{itemize}
\item \Larkita\ always builds the canonical LR(1) automaton; an LALR(1) parser is obtained by merging states
      (see below).  The Earley and CYK parsers of \Lark\ are not available.
\item The parser returns nested tuples.  There are no trees, transformers, or visitors.  Therefore, the rule
      modifiers \verb|?rule| and \verb|!rule| are not supported, and rules and terminals whose names start
      with an underscore are neither inlined nor filtered.
\item Aliases and \verb|%ignore| after an alternative and \verb|%ignore:| control the heads of the tuples.
\item Grammar templates, post-lexers like the \texttt{Indenter}, lexer callbacks, the interactive parser,
      and caching are not supported.
\item Shift/reduce conflicts are resolved as shift with a warning, and every parser comes with an output file
      that describes the LR(1) automaton.
\item By default, the states of the canonical LR(1) automaton are not merged.  If the option
      \verb|merge_states=True| is given, states with the same core are merged as far as this is safe.
\end{itemize}

\end{document}
